\begin{preface}{2026年}{松野裕}
2022年11月、OpenAIがChatGPTを公開したとき、世界は大きく変わった。
わずか2か月で1億人のユーザーを獲得したこの対話型AIは、研究者の日常にも急速に浸透し始めた。
文献調査、論文執筆、プログラミング、データ分析——あらゆる研究活動において、
生成AIは無視できない存在となっている。

しかし、その急速な普及に対して、大学院教育における体系的な指導は追いついていない。
多くの大学院生は、SNSやネット記事の断片的な情報を頼りに、手探りでAIツールを使っている。
ある学生は便利さに感動して過度に依存し、別の学生は「使ってはいけないもの」と
誤解して一切触れようとしない。どちらも、研究者として望ましい姿勢とは言えない。
本書は、この状況を踏まえて執筆したものである。大学院生が生成AIを批判的に、
かつ効果的に活用するための体系的なガイドとして、本書を執筆した。

筆者は2025年度より、日本大学大学院における「情報論」の授業で、
大学院生を対象にAI活用の指導を行ってきた。
授業を通じて痛感したのは、(1)研究活動におけるAI活用に焦点を当てた体系的な教材の不在、
(2)プロンプトエンジニアリングの理論と実践のギャップ、
(3)何をAIに任せてよいのかという倫理的な判断力の必要性、の3点である。
2025年度、2026年度の2年間にわたる授業実践を通じて蓄積した知見と、
受講生からのフィードバックが、本書の土台となっている。

本書を貫く基本哲学は、「批判的AI活用」（Critical AI Usage）である。
これは、AIを盲目的に信頼するのでも、頭ごなしに拒否するのでもない。
AIの出力を常に批判的に検証し、自らの専門知識と判断力をもって取捨選択する姿勢のことである。
具体的には、(1) AIは助手であり、研究者の代替ではない、
(2) AIの出力は常に検証する、(3) AIの使用を透明にする、という3つの原則に基づいている。

本書は、半期15回の授業に対応した構成となっている。第1章から第10章までの各章が
1〜2回の授業に対応し、講義と演習を組み合わせて進めることを想定している。
自習で読み進める場合は、まず第1章と第2章を通読し、基本的な考え方と技法を
身につけることを勧める。その後は、自分の研究活動で最も必要な章から読み進めてよい。

本書は、特定のAIツールに依存しない設計としている。
本文中ではChatGPT、Claude、Geminiなど多様なツールを例として取り上げるが、
個々のツールの操作方法よりも、どのツールにも共通する活用の原則と思考法を重視している。
AI分野の進化は速く、書籍の出版後にも新たなツールや技法が登場することが予想されるため、
GitHubリポジトリにて補足資料や更新情報を提供する予定である。

本書の執筆にあたり、多くの方々にお世話になった。
日本大学大学院理工学研究科の「情報論」受講生の皆さんには、
授業を通じて多くのフィードバックをいただいた。
学生の率直な質問や試行錯誤の経験が、本書の内容を大きく豊かにしてくれた。
心より感謝する。
コロナ社編集部には、企画段階から丁寧なご助言をいただいた。
Webサポートとの連携という出版形態を快くお認めいただいたことに感謝する。
最後に、執筆期間中、常に支えてくれた家族に感謝したい。
\end{preface}
